\documentclass[conference]{IEEEtran}
\IEEEoverridecommandlockouts
\usepackage{amsmath,amssymb,amsfonts}
\usepackage{booktabs}
\usepackage{tabularx}
\usepackage{graphicx}
\usepackage{xcolor}
\usepackage{cite}
\usepackage{microtype}

\title{Mathematical Formulation of the Proactive Golden-Hour Bridge Model (PGBM) for Post-Earthquake UAV-Assisted Response}

\author{\IEEEauthorblockN{Author Name}
\IEEEauthorblockA{\textit{Department of Computer Science and Engineering}\\
\textit{University Name}\\
Email: author@email.com}
}

\begin{document}

\maketitle

\begin{abstract}
This document presents the mathematical formulation of the Proactive Golden-Hour Bridge Model (PGBM).
\end{abstract}

\section{Introduction}
\subsection{Problem in One Sentence}
After Scout UAVs detect survivors, the Command Center must decide at each event time $t$ which available Responder UAVs should execute which feasible tours---ordered sequences $\pi_j=(i_1,\dots,i_k)$---to maximize expected lifesaving value under battery, payload-item, and 3D reachability constraints, before ground rescue arrives.

\section{System Model}
% Define entities, environment, tasks

\section{Problem Formulation}
% Offloading + Routing

\section{Conclusion}
% Summary

\end{document}
